\documentclass[reqno, 12pt]{amsart}
\usepackage[brazil]{babel}
\usepackage[utf8]{inputenc}
\usepackage[T1]{fontenc}
\usepackage{lmodern}
\usepackage{amsfonts,amssymb,amsmath,textcomp,amsthm}
\usepackage{geometry}
\geometry{left=25mm, right=25mm, top=22mm, bottom=28mm}
\usepackage{setspace}
\onehalfspacing
\usepackage{csquotes}
\usepackage[hidelinks]{hyperref}
\usepackage[backend=biber, style=numeric, maxbibnames=99]{biblatex}
\addbibresource{refs_antiramsey.bib}
\setlength{\parindent}{1.2cm}
\numberwithin{equation}{section}

% -------------------------------------------------------
% Ambientes
% -------------------------------------------------------
\theoremstyle{plain}
\newtheorem{theorem}{Teorema}[section]
\newtheorem{proposition}[theorem]{Proposição}
\newtheorem{lemma}[theorem]{Lema}
\newtheorem{corollary}[theorem]{Corolário}

\theoremstyle{definition}
\newtheorem{definition}[theorem]{Definição}
\newtheorem{question}[theorem]{Questão}
\newtheorem{problem}[theorem]{Problema}
\newtheorem{remark}[theorem]{Observação}
\newtheorem{example}[theorem]{Exemplo}

% -------------------------------------------------------
% Macros
% -------------------------------------------------------
\newcommand{\Gnp}{\mathcal{G}(n,p)}
\newcommand{\rb}{\xrightarrow{\,\mathrm{rb}\,}}
\newcommand{\ramsey}[1]{\xrightarrow{\,#1\,}}
\newcommand{\cram}{\xrightarrow{\,\mathrm{c\text{-}ram}\,}}
\newcommand{\mdens}{m_2}
\newcommand{\maxdens}{m}
\newcommand{\whp}{a.a.s.}

\date{\today}

% ============================================================
\begin{document}
% ============================================================

\begin{center}
  {\large\textbf{Ramsey Assimétrico e Anti-Ramsey em Grafos Aleatórios}}

  \vspace{0.4cm}

  Afonso Sant'Anna

  \vspace{0.2cm}

  {\small Orientador: Yoshiharu Kohayakawa \quad
   Coorientador: Guilherme O.\ Mota}

  \vspace{0.2cm}

  {\small Instituto de Matemática e Estatística -- USP}
\end{center}

\vspace{0.5cm}

\noindent\textbf{Nota.}
Este documento é um rascunho de trabalho para o doutorado de Afonso Sant'Anna.
Ele reúne definições, teoremas e resultados dos artigos da literatura relevante,
organizados para facilitar a escrita dos próprios resultados do aluno.
Seções marcadas com \emph{(RA)} referem-se ao tema de Ramsey Assimétrico/Canônico;
seções marcadas com \emph{(AR)} referem-se ao tema Anti-Ramsey.

\tableofcontents

\bigskip

% ============================================================
\section{Preliminares e Notação}
% ============================================================

Denotamos por $\Gnp$ o grafo aleatório binomial com $n$ vértices em que cada aresta
está presente independentemente com probabilidade $p = p(n)$.
Dizemos que $\Gnp$ satisfaz uma propriedade \emph{assintoticamente quase certamente}
(\whp) se a probabilidade de satisfazê-la tende a $1$ quando $n \to \infty$.

\begin{definition}[Limiar e tipos de limiar]
Seja $\mathcal{P}$ uma propriedade monótona crescente de grafos.
Uma função $\hat{p} = \hat{p}(n)$ é um \textbf{limiar grosseiro} (\emph{coarse threshold})
para $\mathcal{P}$ em $\Gnp$ se
\[
  \lim_{n\to\infty} \Pr[\Gnp \in \mathcal{P}] =
  \begin{cases}
    0 & \text{se } p = o(\hat{p}),\\
    1 & \text{se } p = \omega(\hat{p}).
  \end{cases}
\]
Dizemos que $\hat{p}$ é um \textbf{limiar semi-agudo} (\emph{semi-sharp threshold})
se existem constantes $c, C > 0$ tais que
\[
  \lim_{n\to\infty} \Pr[\Gnp \in \mathcal{P}] =
  \begin{cases}
    0 & \text{se } p \le c\,\hat{p},\\
    1 & \text{se } p \ge C\,\hat{p}.
  \end{cases}
\]
O \textbf{$0$-enunciado} refere-se à afirmação para $p$ abaixo do limiar,
e o \textbf{$1$-enunciado} à afirmação para $p$ acima do limiar.
\end{definition}

\begin{definition}[Densidades de um grafo]
\label{def:densidades}
Dado um grafo $H$, definem-se:
\begin{itemize}
  \item A \textbf{densidade} de $H$:
    $d(H) := \dfrac{e(H)}{v(H)}$, onde $e(H) = |E(H)|$ e $v(H) = |V(H)|$.

  \item A \textbf{densidade máxima} de $H$:
    \[
      \maxdens(H) := \max\!\left\{\frac{e(J)}{v(J)} : J \subseteq H,\; v(J) \ge 1\right\}.
    \]
    Dizemos que $H$ é \textbf{balanceado} se $\maxdens(H) = d(H)$, e
    \textbf{estritamente balanceado} se $\maxdens(H) > d(J)$ para todo
    $J \subsetneq H$ com $v(J) \ge 1$.

  \item A \textbf{2-densidade} de $H$ (para $v(H) \ge 3$):
    \[
      d_2(H) := \frac{e(H)-1}{v(H)-2}.
    \]

  \item A \textbf{2-densidade máxima} de $H$:
    \[
      \mdens(H) := \max\!\left\{d_2(J) : J \subseteq H,\; v(J) \ge 3\right\}
        \cup \left\{\tfrac{1}{2}\right\}.
    \]
    (A fração $\frac{1}{2}$ corresponde ao caso $H \cong K_2$.)
    Dizemos que $H$ é \textbf{$2$-balanceado} se $\mdens(H) = d_2(H)$, e
    \textbf{estritamente $2$-balanceado} se $\mdens(H) > d_2(J)$ para todo
    $J \subsetneq H$.
\end{itemize}
\end{definition}

\begin{example}
  Para o clique $K_k$ com $k \ge 3$:
  $\mdens(K_k) = \dfrac{\binom{k}{2}-1}{k-2} = \dfrac{k+1}{2}$.
  Para o ciclo $C_\ell$ com $\ell \ge 3$:
  $\mdens(C_\ell) = \dfrac{\ell-1}{\ell-2}$.
  Em particular,
  $\mdens(K_5) = 3$, $\mdens(C_5) = \frac{4}{3}$, $\mdens(C_4) = \frac{3}{2}$.
\end{example}

% ============================================================
\section{Propriedade de Ramsey em Grafos Aleatórios}
% ============================================================

\begin{definition}[Propriedade $r$-Ramsey]
Para grafos $G$ e $H$ e um inteiro $r \ge 2$, escrevemos
$G \ramsey{r} H$ para denotar a propriedade de que,
em toda coloração das arestas de $G$ com $r$ cores,
existe pelo menos uma cópia monocromática de $H$ em $G$.
\end{definition}

O seguinte resultado é o marco fundamental da teoria:

\begin{theorem}[Rödl--Ruciński \cite{rodl1993lower, rodl1994random, rodl1995threshold}]
\label{thm:rodl-rucinski}
Seja $r \ge 2$ um inteiro e seja $H$ um grafo que não é uma floresta de estrelas
(e, no caso $r = 2$, também não é um caminho de comprimento~$3$).
Existem constantes positivas $c = c(H,r)$ e $C = C(H,r)$ tais que
\[
  \lim_{n\to\infty} \Pr\!\bigl[\Gnp \ramsey{r} H\bigr] =
  \begin{cases}
    0 & \text{se } p \le c\, n^{-1/\mdens(H)},\\
    1 & \text{se } p \ge C\, n^{-1/\mdens(H)}.
  \end{cases}
\]
Em particular, $\hat{p}(n) = n^{-1/\mdens(H)}$ é o limiar semi-agudo para
a propriedade $r$-Ramsey de $H$ em $\Gnp$.
\end{theorem}

Uma prova curta e elegante do Teorema~\ref{thm:rodl-rucinski} foi dada por
Nenadov e Steger \cite{nenadov2016short}, usando a técnica de
\emph{containers} de hipergrafos.

% ============================================================
\section{Propriedade Anti-Ramsey em Grafos Aleatórios} \label{sec:antiramsey}
% ============================================================

\begin{definition}[Propriedade anti-Ramsey]
\label{def:antiramsey}
Dados grafos $G$ e $H$, escrevemos $G \rb H$ para denotar a
\textbf{propriedade anti-Ramsey} de $G$ com respeito a $H$:
para toda coloração própria das arestas de $G$
(com número arbitrário de cores),
existe pelo menos uma cópia \emph{arco-íris} de $H$ em $G$,
isto é, uma cópia de $H$ em que cada aresta recebe uma cor distinta.
\end{definition}

\begin{remark}
A propriedade $G \rb H$ é monótona crescente em $G$.
Logo, para qualquer grafo fixo $H$, o grafo aleatório $\Gnp$ admite
um limiar $p_H^{\mathrm{rb}} = p_H^{\mathrm{rb}}(n)$ para essa propriedade
(cf.~Bollobás--Thomason).
\end{remark}

A questão central é determinar $p_H^{\mathrm{rb}}$ em termos de parâmetros
combinatórios de $H$.

\subsection{1-Enunciado: limite superior para o limiar}

O primeiro resultado geral foi obtido por Rödl e Tuza~\cite{rodl1992rainbow}
para ciclos, e depois estendido por Kohayakawa, Konstadinidis e Mota:

\begin{theorem}[Kohayakawa--Konstadinidis--Mota \cite{kohayakawa2014anti}]
\label{thm:1statement}
Seja $H$ um grafo qualquer. Existe uma constante $C > 0$ tal que,
se $p = p(n) \ge C\, n^{-1/\mdens(H)}$, então \whp\
$\Gnp \rb H$.
Em particular, $p_H^{\mathrm{rb}} \le n^{-1/\mdens(H)}$.
\end{theorem}

Em outras palavras, o limiar anti-Ramsey nunca está acima do limiar Ramsey clássico.

\subsection{0-Enunciado: limite inferior — o \emph{framework} de Nenadov et al.}

Para o limite inferior, Nenadov, Person, Škorić e Steger desenvolveram um
arcabouço geral que reduz o problema probabilístico a uma questão determinística:

\begin{theorem}[Nenadov--Person--Škorić--Steger \cite{nenadov2017algorithmic}]
\label{thm:framework}
Seja $H$ um grafo estritamente $2$-balanceado.
Se todo grafo $G$ tal que $G \rb H$ satisfaz $\maxdens(G) > \mdens(H)$,
então existe uma constante $c > 0$ tal que, se $p \le c\, n^{-1/\mdens(H)}$,
então \whp\ $\Gnp \not\rb H$.

Em outras palavras, verificar a condição determinística acima
é suficiente para obter o $0$-enunciado semi-agudo.
\end{theorem}

\begin{definition}[Condição determinística anti-Ramsey]
\label{def:cond-det}
Dado um grafo $H$, dizemos que $H$ \textbf{satisfaz a condição}~(D)
se todo grafo $G$ tal que $G \rb H$ satisfaz $\maxdens(G) > \mdens(H)$.
\end{definition}

\subsection{Resultados para classes específicas de grafos}

\subsubsection*{Ciclos}

\begin{theorem}[Barros--Cavalar--Mota--Parczyk \cite{AntiCycles}]
\label{thm:anticycles}
O limiar $p_{C_\ell}^{\mathrm{rb}}$ para a propriedade anti-Ramsey do ciclo $C_\ell$
satisfaz:
\begin{itemize}
  \item Se $\ell \ge 5$, então
    $p_{C_\ell}^{\mathrm{rb}} = n^{-1/\mdens(C_\ell)} = n^{-(\ell-2)/(\ell-1)}$
    (limiar semi-agudo).
  \item Para $\ell = 4$:
    $p_{C_4}^{\mathrm{rb}} = n^{-3/4} \ll n^{-1/\mdens(C_4)} = n^{-2/3}$.
  \item Para $\ell = 3$: a condição de ser grafo arco-íris é satisfeita
    trivialmente, logo
    $p_{C_3}^{\mathrm{rb}} = n^{-1}$ (limiar para a presença de um triângulo).
\end{itemize}
\end{theorem}

\subsubsection*{Grafos completos}

\begin{theorem}[Kohayakawa--Mota--Parczyk--Schnitzer \cite{AntiComplete}]
\label{thm:anticomplete}
Para o clique $K_k$, o limiar $p_{K_k}^{\mathrm{rb}}$ satisfaz:
\begin{itemize}
  \item Se $k \ge 5$, então
    $p_{K_k}^{\mathrm{rb}} = n^{-1/\mdens(K_k)} = n^{-2/(k+1)}$
    (limiar semi-agudo).
  \item Para $k = 4$:
    $p_{K_4}^{\mathrm{rb}} = n^{-7/15} \ll n^{-1/\mdens(K_4)} = n^{-2/5}$.
  \item Para $k = 3$: trivialmente $p_{K_3}^{\mathrm{rb}} = n^{-1}$.
\end{itemize}
\end{theorem}

\subsubsection*{Caso geral: 2-densidade alta}

\begin{theorem}[Kuperwasser \cite{kuperwasser2026anti}]
\label{thm:kuperwasser}
Seja $H$ um grafo com $\mdens(H) \ge 19$.
Existem constantes $c_0 < c_1$ tais que
\[
  \lim_{n\to\infty}\Pr\!\bigl[\Gnp \rb H\bigr] =
  \begin{cases}
    0 & \text{se } p \le c_0\, n^{-1/\mdens(H)},\\
    1 & \text{se } p \ge c_1\, n^{-1/\mdens(H)}.
  \end{cases}
\]
\end{theorem}

O Teorema~\ref{thm:kuperwasser} é provado verificando a condição determinística
do Teorema~\ref{thm:framework} através do seguinte resultado:

\begin{theorem}[Kuperwasser \cite{kuperwasser2026anti}]
\label{thm:kuperwasser-det}
Todo grafo $G$ com $\maxdens(G) \ge 18$ admite uma coloração própria das arestas
tal que todo subgrafo arco-íris tem $2$-densidade estritamente menor que $\maxdens(G)$.
Em particular, se $\mdens(H) > 18$, então $H$ satisfaz a condição~\emph{(D)}.
\end{theorem}

\subsection{A tabela de limiares conhecidos}

A tabela a seguir resume o estado atual dos limiares para a propriedade anti-Ramsey,
adaptada de \cite{behague2025thresholds}:

\begin{center}
\renewcommand{\arraystretch}{1.4}
\begin{tabular}{lll}
\hline
\textbf{Grafo $H$} & \textbf{Limiar $p_H^{\mathrm{rb}}$} & \textbf{Referência}\\
\hline
$K_\ell$, $\ell \ge 5$ & $n^{-2/(\ell+1)}$ (semi-agudo) & \cite{AntiComplete}\\
$K_4$                  & $n^{-7/15} \ll n^{-1/m_2(K_4)}$& \cite{AntiComplete}\\
$K_3$                  & $n^{-1}$                        & trivial\\
$C_\ell$, $\ell \ge 5$ & $n^{-(\ell-2)/(\ell-1)}$ (semi-agudo) & \cite{AntiCycles}\\
$C_4$                  & $n^{-3/4} \ll n^{-1/m_2(C_4)}$ & \cite{AntiCycles}\\
$C_3$                  & $n^{-1}$                        & trivial\\
$H$ com $m_2(H)\ge 19$ & $n^{-1/m_2(H)}$ (semi-agudo)   & \cite{kuperwasser2026anti}\\
\hline
\end{tabular}
\end{center}

\subsection{Questão em aberto: grafos de cintura grande}

Os resultados acima revelam um padrão: o limiar anti-Ramsey diverge de
$n^{-1/\mdens(H)}$ exatamente para grafos ``pequenos'' de cintura baixa
($K_3$, $K_4$, $C_4$). Essa observação motiva a seguinte questão central
deste projeto de doutorado:

\begin{question}
\label{q:cintura}
Seja $H$ um grafo com cintura (comprimento do menor ciclo) estritamente maior que~$4$.
É verdade que $p_H^{\mathrm{rb}} = n^{-1/\mdens(H)}$?
Equivalentemente, $H$ sempre satisfaz a condição~\emph{(D)}
(Definição~\ref{def:cond-det}) quando $\mathrm{girth}(H) > 4$?
\end{question}

% ============================================================
\section{Propriedade Ramsey Restrita (Constrained Ramsey)} \label{sec:cram}
% ============================================================

Este problema, estudado em \cite{behague2025thresholds}, generaliza tanto
a propriedade de Ramsey clássica quanto a anti-Ramsey.

\begin{definition}[Propriedade Ramsey restrita]
\label{def:cram}
Dados grafos $H_1$ e $H_2$, escrevemos $G \cram (H_1, H_2)$ para denotar
a \textbf{propriedade Ramsey restrita} (\emph{constrained Ramsey}):
para toda coloração própria das arestas de $G$,
existe uma cópia monocromática de $H_1$ \emph{ou} uma cópia arco-íris de $H_2$.
\end{definition}

\begin{remark}
Quando $H_2$ é uma floresta, qualquer coloração própria tem trivialmente uma
cópia arco-íris de $H_2$ (desde que o grafo seja grande o suficiente). Logo,
o caso interessante é quando $H_2$ não é uma floresta.

Quando $H_1 = $ estrela com $k$ arestas $K_{1,k}$ e $H_2 = H$, a propriedade
$G \cram (K_{1,k}, H)$ equivale a dizer que toda coloração \emph{localmente
$(k-1)$-limitada} de $G$ (em que cada cor aparece no máximo $k-1$ vezes por vértice)
contém uma cópia arco-íris de $H$. Assim,
a propriedade anti-Ramsey é o caso particular $k = 2$
(colorações próprias = colorações localmente $1$-limitadas).
\end{remark}

\subsection{Resultados principais de Behague et al.}

\begin{theorem}[Behague--Hancock--Hyde--Letzter--Morrison \cite{behague2025thresholds}]
\label{thm:behague-1-stmt}
Sejam $k \ge 2$ e $H$ um grafo qualquer.
Existe uma constante $C > 0$ tal que, se $p \ge C\, n^{-1/\mdens(H)}$, então
\whp\ $\Gnp \cram (K_{1,k}, H)$.
\end{theorem}

O Teorema~\ref{thm:behague-1-stmt} generaliza o Teorema~\ref{thm:1statement}.

\begin{definition}[Densidades relevantes para Ramsey restrito]
Para o problema restrito com $H_1 = K_{1,k}$, definem-se:
\begin{align*}
  m_a(H) &:= \inf\!\bigl\{\maxdens(G) : G \cram (K_{1,k}, H)\bigr\}\\
         &\phantom{:}= \text{densidade máxima do menor grafo} \ G\ \text{que força }\\
         &\phantom{= \text{densidade máxima do menor grafo} G \text{que força}} \text{Ramsey restrito para } (K_{1,k}, H).
\end{align*}
O limiar candidato natural é então $n^{-1/m_a(H)}$.
\end{definition}

\begin{theorem}[Behague et al.\ \cite{behague2025thresholds} -- $0$-enunciado geral]
\label{thm:behague-0-stmt}
Sejam $k \ge 3$ e $H$ um grafo que não é floresta.
\begin{enumerate}
  \item[(a)] Se $k \ge 4$, \emph{ou} $\mdens(H) \ne 2$, \emph{ou}
    $\mdens(H) = 2$ e $H$ contém um subgrafo estritamente $2$-balanceado
    $J \ne K_3$ com $\mdens(J) = 2$,
    então existem $c, C > 0$ tais que
    \[
      \lim_{n\to\infty} \Pr\!\bigl[\Gnp \cram (K_{1,k}, H)\bigr] =
      \begin{cases}
        0 & \text{se } p \le c\, n^{-1/\mdens(H)},\\
        1 & \text{se } p \ge C\, n^{-1/\mdens(H)}.
      \end{cases}
    \]
  \item[(b)] Se $k = 3$ e $\mdens(H) = 2$ e o único subgrafo estritamente
    $2$-balanceado de $H$ com $2$-densidade~$2$ é $K_3$ (e.g.\ $H = K_3$),
    então o limiar é mais fraco: coarse threshold em $n^{-1/2}$.
\end{enumerate}
\end{theorem}

% ============================================================
\section{Propriedade Ramsey Canônica Assimétrica} \label{sec:canonico}
% ============================================================

A teoria canônica de Ramsey estuda estruturas que emergem em colorações
\emph{arbitrárias} (sem restrição de número de cores).

\begin{definition}[Coloração canônica]
\label{def:canonica}
Seja $c : E(G) \to \mathbb{N}$ uma coloração arbitrária das arestas de $G$,
com uma ordenação total $v_1 < v_2 < \cdots < v_n$ dos vértices de $G$.
Um subgrafo $H \subseteq G$ possui uma \textbf{coloração canônica} se
satisfaz um dos seguintes padrões:
\begin{enumerate}
  \item \textbf{Monocromático:} todas as arestas de $H$ têm a mesma cor.
  \item \textbf{Arco-íris:} todas as arestas de $H$ têm cores distintas.
  \item \textbf{Lexicográfico (min-lex):} a cor de $v_iv_j$ (com $v_i < v_j$)
    depende unicamente de $v_i$.
  \item \textbf{Lexicográfico (max-lex):} a cor de $v_iv_j$ (com $v_i < v_j$)
    depende unicamente de $v_j$.
\end{enumerate}
\end{definition}

\begin{theorem}[Erdős--Rado \cite{erdos1950combinatorial}]
Para qualquer inteiro $\ell \ge 1$, existe $N = N(\ell)$ tal que toda coloração
arbitrária das arestas de $K_N$ contém uma cópia canônica de $K_\ell$.
\end{theorem}

O resultado de Erdős--Rado foi trazido para o contexto aleatório por Kamčev e Schacht:

\begin{theorem}[Kamčev--Schacht \cite{kamcev2023canonical}]
\label{thm:kamcev}
Para todo $\ell \ge 4$, o limiar para a propriedade de que toda coloração
arbitrária de $\Gnp$ contenha uma cópia canônica de $K_\ell$ é
$\hat{p}(n) = n^{-2/(\ell+1)} = n^{-1/\mdens(K_\ell)}$.
\end{theorem}

\subsection{Cenário assimétrico}

\begin{definition}[Propriedade Ramsey canônica assimétrica]
\label{def:assimetrico}
Dados grafos $H_1$, $H_2$ e $H_3$, escrevemos
\[
  G \xrightarrow{\mathrm{can}} (H_1, H_2, H_3)
\]
para denotar a seguinte propriedade: para toda coloração arbitrária das
arestas de $G$, existe pelo menos
\begin{itemize}
  \item uma cópia \textbf{monocromática} de $H_1$, ou
  \item uma cópia \textbf{arco-íris} de $H_2$, ou
  \item uma cópia \textbf{lexicográfica} de $H_3$.
\end{itemize}
\end{definition}

\begin{question}
\label{q:assimetrico}
Qual é o comportamento da função limiar $\hat{p}(n)$ para a propriedade
$\Gnp \xrightarrow{\mathrm{can}} (H_1, H_2, H_3)$
quando $H_1$, $H_2$ e $H_3$ não são todos iguais?
Em particular, como o limiar depende das densidades $\mdens(H_i)$?
\end{question}

% ============================================================
\section{O Caso Assimétrico \texorpdfstring{$(K_3, K_4, K_5)$}{(K3, K4, K5)}}
% ============================================================

Queremos determinar o limiar exato para a propriedade $\Gnp \xrightarrow{\mathrm{can}} (K_3,K_4,K_5)$.
Note que, para $p \ll n^{-1/\mdens(K_3)} = n^{-1/2}$, o grafo aleatório \whp\ não possui triângulos, tornando a propriedade trivialmente falsa (através de uma coloração arco-íris, por exemplo).
Por outro lado, para $p \gg n^{-1/\mdens(K_5)} = n^{-1/3}$, a existência de um $K_5$ canônico é garantida \whp, o que força a existência de uma das estruturas proibidas.

O resultado principal que buscamos formalizar é que o limiar coincide com o do clique mais denso.
Para a prova do $0$-enunciado, precisaremos das seguintes definições estruturais:

\begin{definition}[$H$-fechamento]
\label{def:H-fechado}
Dado um grafo $G$ e um grafo $H$, dizemos que:
\begin{itemize}
  \item Uma aresta $e \in E(G)$ é \textbf{$H$-fechada} se pertence a pelo menos duas cópias de $H$ em $G$.
  \item Uma cópia $H' \subseteq G$ de $H$ é \textbf{$H$-fechada} se pelo menos três de suas arestas são $H$-fechadas.
  \item O grafo $G$ é \textbf{$H$-fechado} se todo vértice e toda aresta de $G$ pertencem a pelo menos uma cópia de $H$, e toda cópia de $H$ em $G$ é $H$-fechada.
\end{itemize}
\end{definition}

\begin{definition}[$H$-Bloco]
\label{def:H-bloco}
Um grafo $G$ é um \textbf{$H$-bloco} se $G$ é $H$-fechado e, para todo subconjunto próprio não-vazio de arestas $E' \subsetneq E(G)$, existe uma cópia $H'$ de $H$ em $G$ tal que $E(H') \cap E' \ne \emptyset$ e $E(H') \setminus E' \ne \emptyset$.
\end{definition}

\begin{remark}
A chave do framework de Nenadov et al.\ é que, quando $p \ll n^{-1/\maxdens(H)}$, os $H$-blocos presentes em $\Gnp$ têm tamanho $O(1)$ \whp.
\end{remark}

\begin{theorem}
  Existe um limiar semi-agudo para a propriedade $\Gnp \xrightarrow{\mathrm{can}} (K_3,K_4,K_5)$ em $\hat{p}(n) = n^{-1/3}$. Ou seja, existem constantes $c, C > 0$ tais que:
  \[
  \lim_{n \to \infty} \mathbb{P}(\Gnp \xrightarrow{\mathrm{can}} (K_3,K_4,K_5)) = 
  \begin{cases}
    1 & \text{se } p \ge C n^{-1/3}, \\
    0 & \text{se } p \le c n^{-1/3}.  
  \end{cases}
  \]
\end{theorem}

\begin{proof}[Esbo\c{c}o da Prova]
  O $1$-enunciado segue do Teorema de Kam\v{c}ev--Schacht~\cite{kamcev2023canonical}.

  Para o $0$-enunciado, descrevemos o algoritmo \textsc{Colora\c{c}\~{a}oCan} que, dado $G$ com $\maxdens(G) \le 3$, constr\'{o}i explicitamente uma colora\c{c}\~{a}o das arestas de $G$ sem $K_3$ monocrom\'{a}tico, $K_4$ arco-\'{i}ris ou $K_5$ lexicogr\'{a}fico.

  \textbf{Passo 1 --- Colora\c{c}\~{a}o Lexicogr\'{a}fica Inicial.}
  Fixamos uma ordena\c{c}\~{a}o total $\sigma: V(G) \to [n]$ e uma atribui\c{c}\~{a}o $\varphi: V(G) \to \mathbb{N}$.
  Definimos a colora\c{c}\~{a}o $c: E(G) \to \mathbb{N}$ por
  \[ c(uv) \;:=\; \varphi\!\bigl(\min_\sigma\{u,v\}\bigr), \]
  ou seja, a cor de cada aresta \'{e} determinada pelo v\'{e}rtice de menor \'{i}ndice segundo $\sigma$.
  Por constru\c{c}\~{a}o, $c$ \'{e} uma colora\c{c}\~{a}o min-lex e satisfaz:
  \begin{itemize}
    \item n\~{a}o existe $K_3$ monocrom\'{a}tico sob $c$;
    \item n\~{a}o existe $K_4$ arco-\'{i}ris sob $c$.
  \end{itemize}
  O problema \'{e} que toda c\'{o}pia de $K_5$ em $G$ \'{e} lexicogr\'{a}fica (min-lex) sob $c$ --- exatamente a estrutura proibida que precisamos eliminar.

  \textbf{Passo 2 --- Remo\c{c}\~{a}o de Arestas Livres.}
  Removemos de $G$ todas as arestas que n\~{a}o pertencem a nenhuma c\'{o}pia de $K_5$.
  Seja $G' \subseteq G$ o grafo resultante. As arestas removidas mant\^{e}m a colora\c{c}\~{a}o do Passo~1 e n\~{a}o causam problemas.

  \textbf{Passo 3 --- Decomposi\c{c}\~{a}o em $K_5$-Blocos.}
  Decompomos $G'$ em seus $K_5$-blocos (Defini\c{c}\~{a}o~\ref{def:H-bloco}).
  Como $\maxdens(G) \le \mdens(K_5) = 3$, cada $K_5$-bloco tem tamanho $O(1)$.

  \textbf{Passo 4 --- Recolora\c{c}\~{a}o dos Blocos.}
  Para cada $K_5$-bloco $B$, substitu\'{i}mos a colora\c{c}\~{a}o $c|_{E(B)}$ por uma nova colora\c{c}\~{a}o que evita $K_5$ lexicogr\'{a}fico, sem introduzir $K_3$ monocrom\'{a}tico ou $K_4$ arco-\'{i}ris.
  A exist\^{e}ncia de tal colora\c{c}\~{a}o em cada bloco \'{e} garantida pelo Lema~\ref{lem:0-statement-K5} abaixo.
\end{proof}

\begin{lemma}
  \label{lem:0-statement-K5}
  Seja $B$ um $K_5$-bloco com $\maxdens(B) \le 3$. Ent\~{a}o existe uma colora\c{c}\~{a}o $c: E(B) \to \mathbb{N}$ tal que:
  \begin{enumerate}
    \item[(i)] $c$ n\~{a}o possui $K_3$ monocrom\'{a}tico;
    \item[(ii)] $c$ n\~{a}o possui $K_4$ arco-\'{i}ris;
    \item[(iii)] $c$ n\~{a}o possui $K_5$ lexicogr\'{a}fico.
  \end{enumerate}
\end{lemma}

\begin{proof}
  
\end{proof}

\vspace{0.5cm}

% ============================================================
\section{Resultados em Andamento}
% ============================================================

\emph{[Esta seção é para registrar resultados e progressos do aluno.]}

\vspace{1cm}

% ============================================================
\printbibliography
% ============================================================

\end{document}
